\documentclass[10pt]{ctexart}
\usepackage[a4paper,margin=0.9in]{geometry}
\usepackage{amsmath,amssymb,booktabs,array,tabularx}
\usepackage{enumitem}
\usepackage{listings}
\usepackage{xcolor}
\usepackage{microtype}
\usepackage{hyperref}
\usepackage{xurl}

\setlist[itemize]{leftmargin=1.5em}
\setlist[enumerate]{leftmargin=1.8em}
\setlength{\parskip}{0.35em}
\setlength{\parindent}{0em}
\renewcommand{\arraystretch}{1.15}

\lstset{
  basicstyle=\ttfamily\small,
  breaklines=true,
  columns=fullflexible,
  keepspaces=true,
  frame=single,
  framerule=0.2pt,
  xleftmargin=0.5em,
  xrightmargin=0.5em
}

\hypersetup{
  colorlinks=true,
  linkcolor=blue,
  citecolor=blue,
  urlcolor=blue
}

\title{022：Self-Summary Credit Assignment：\\用轨迹自总结连接长程 Agent RL 的过程监督与策略改进}
\author{}
\date{2026-06-26}

\begin{document}
\maketitle
\vspace{-1.2em}

\section*{摘要}
本文将一个新的长程 Agent RL 训练思路整理为可落地实验方案：在 rollout 中显式插入由 actor 自己生成的 summary / policy-improvement note，形成
\begin{lstlisting}
prompt -> traj1 -> summary -> traj2 -> reward
\end{lstlisting}

其中 summary 不是旁路日志，而是一个可训练的中间动作。它一方面被后续 \path{traj2} 的 reward 训练，另一方面可以作为 hindsight condition 反过来评价 \path{traj1} 中哪些 step / token / action 更值得强化或抑制。本文暂称该方向为 \textbf{Self-Summary Credit Assignment, SSCA}。

本文定位不是替代 GiGPO / GAGPO / StepPPO，而是补充一个新的 credit source：当 exact same-state grouping 稀疏、local rerollout 昂贵、critic 学习慢时，利用“summary-conditioned policy shift”给 prefix trajectory 提供低成本的隐式过程监督。

\section{核心问题}
当前长程 Agent RL 的主要困难不是只有 sparse reward，而是三件事同时出现：
\begin{enumerate}
  \item rollout 昂贵，不能为每个中间 state 大量重采样；
  \item trajectory 很长，episode reward 直接广播会污染大量 routine steps；
  \item state 复现稀疏，GiGPO / GAGPO 依赖的 same-state group 在开放式环境中不稳定。
\end{enumerate}

已有方案大致分成三类：
\begin{itemize}
  \item \textbf{Within-batch step credit}：GiGPO / GAGPO 用同状态或 grouped value proxy 从当前 batch 中构造 step-level advantage。
  \item \textbf{Reflection / retry training}：Reflexion、Self-Refine、Reflect, Retry, Reward 证明反思可以改善下一次尝试。
  \item \textbf{Policy-implied dense credit}：VIMPO、SC-GRPO、SDPO 用 policy log-ratio 或 self-conditioned distribution shift 提供 token-level credit。
\end{itemize}

SSCA 的切口是把这三条线合并：
\begin{quote}
summary 不只是给后续决策用的上下文，也可以作为 prefix trajectory 的 hindsight policy condition，从而构造 prefix-level / step-level / token-level 的隐式 advantage 和 value target。
\end{quote}

\section{Rollout 结构}
对同一个任务 prompt $x$，将一条 trajectory 切成三段：
\begin{equation}
  \tau = (\tau_1,\sigma,\tau_2).
\end{equation}

其中：
\begin{itemize}
  \item $\tau_1=(o_0,a_0,\ldots,o_k,a_k)$：前半段尝试；
  \item $r_1$：阶段性 feedback，可来自 verifier、环境中间分、失败原因、工具 stderr、GUI state diff 或测试结果；
  \item $\sigma=\mathrm{Summary}_\theta(\tau_1,r_1)$：actor 生成的总结、失败归因和下一步策略改进；
  \item $\tau_2\sim\pi_\theta(\cdot\mid x,\tau_1,\sigma)$：actor 在 summary 条件下继续执行；
  \item $R=R(\tau_1,\sigma,\tau_2)$：suffix reward 或 final reward。
\end{itemize}

建议第一版 summary prompt 固定成结构化格式：
\begin{lstlisting}
Summarize the previous trajectory using the observed feedback.
1. State belief: what is currently known?
2. Failure / success causes: which decisions helped or hurt?
3. Policy improvement: what should be done differently next?
4. Next-step plan: what should the agent do in the continuation?
\end{lstlisting}

这样做的目的不是追求漂亮自然语言，而是让 summary 的信息密度、可验证性和可训练性更稳定。

\section{训练目标}
\subsection{目标一：用 suffix reward 训练 summary 和 traj2}
最直接的目标是让 summary 和后续动作共同解释 suffix improvement。定义 prefix-conditioned advantage：
\begin{equation}
  A_{\mathrm{suffix}} = R(\tau_1,\sigma,\tau_2) - b(x,\tau_1).
\end{equation}

这里 baseline 不应只用同 prompt group mean。更合理的是估计“给定同一个 prefix，不靠当前 summary 时的 continuation 期望”：
\begin{equation}
  b(x,\tau_1) \approx \mathbb{E}[R\mid x,\tau_1,\mathrm{no\ summary/generic\ summary}].
\end{equation}

第一版可以分三档实现：
\begin{center}
\begin{tabularx}{\linewidth}{lllX}
\toprule
版本 & baseline & 优点 & 风险 \\
\midrule
V0 & prompt group mean & 最简单，便于 smoke test & prefix quality 和 summary quality 混在一起 \\
V1 & 同一 $\tau_1$ 下多条 continuation & 因果更干净 & 需要额外 rerollout \\
V2 & $V_\phi(x,\tau_1)$ & rollout 成本低 & value 早期不稳 \\
\bottomrule
\end{tabularx}
\end{center}

suffix loss 可写为：
\begin{equation}
  \mathcal{L}_{\mathrm{suffix}}
  = - \sum_{t\in\sigma\cup\tau_2}
      \mathrm{clip\_ratio}_t\cdot \mathrm{sg}(A_{\mathrm{suffix}}).
\end{equation}

为了避免空泛 summary 被 final reward 顺带奖励，建议加 gate：
\begin{equation}
  g_{\mathrm{useful}}=\mathbf{1}[R(\tau_2)>b(x,\tau_1)]
\end{equation}
或更严格地只奖励“$\tau_1$ 失败但 $\tau_2$ 成功”的 summary。这与 Reflect, Retry, Reward 的经验一致：只有能把 retry 变好的 reflection 才值得奖励。

\subsection{目标二：用 summary-conditioned shift 监督 traj1}
VIMPO 用 policy-reference log-ratio 构造 policy-implied value / advantage：
\begin{equation}
  \rho_t = \beta \log\frac{\pi_\theta(a_t\mid s_t)}{\pi_{\mathrm{ref}}(a_t\mid s_t)}.
\end{equation}

SSCA 的关键替换是：不把分母设为固定 reference，而是构造同一模型在两个上下文下的 counterfactual policy：
\begin{align}
  \pi_{\mathrm{now}}(a_t\mid h_t) &= \pi_\theta(a_t\mid x,h_t),\\
  \pi_{\mathrm{sum}}(a_t\mid h_t,\sigma) &= \pi_\theta(a_t\mid x,\sigma,h_t).
\end{align}

其中 $h_t$ 是 \path{traj1} 的 prefix。定义 summary-conditioned log-ratio：
\begin{equation}
  \Delta_t^{\mathrm{sum}}
  = \log\pi_{\mathrm{sum}}(a_t\mid h_t,\sigma)
    - \log\pi_{\mathrm{now}}(a_t\mid h_t).
\end{equation}

直觉是：如果模型看到 summary 后更认可 \path{traj1} 的某个动作，则 $\Delta_t^{\mathrm{sum}}>0$；如果 summary 暗示该动作是错误分叉，则 $\Delta_t^{\mathrm{sum}}<0$。这让 summary 成为 \path{traj1} 的 hindsight evaluator。

为了降低自举风险，第一版不要让 numerator / denominator 同时反传。建议采用 stop-gradient weight：
\begin{equation}
  \mathcal{L}_{\mathrm{prefix}}
  = -\sum_{t\in\tau_1}
    \mathrm{sg}\left(g_{\mathrm{useful}}\cdot
    \mathrm{Normalize}(\Delta_t^{\mathrm{sum}})\right)
    \log \pi_\theta(a_t\mid h_t).
\end{equation}

也可以加入 KL-centering 的 VIMPO-style 版本：
\begin{equation}
  A_t^{\mathrm{sum}}
  = \beta\Delta_t^{\mathrm{sum}}
  - \beta\,\mathrm{KL}\!\left(\pi_{\mathrm{sum}}(\cdot\mid h_t,\sigma)
    \Vert \pi_{\mathrm{now}}(\cdot\mid h_t)\right),
\end{equation}
但第一版可以先用 sampled-token score、top-k KL 或 candidate-set KL 近似，避免完整 vocab KL 的额外开销。

\subsection{目标三：用 summary 改善 value prediction}
summary 还可以作为 value predictor 的外部信息源。最直接的 learned value 版本是：
\begin{align}
  V_\phi(h_t) &\rightarrow \text{summary 前 expected return},\\
  V_\phi(h_t,\sigma) &\rightarrow \text{summary 后 expected return},\\
  \Delta V_\phi(h_t,\sigma) &= V_\phi(h_t,\sigma)-V_\phi(h_t).
\end{align}

如果 summary 有用，应满足：
\begin{equation}
  \Delta V_\phi(h_k,\sigma)
  \approx R(\tau_2)-b(x,\tau_1).
\end{equation}

也可以构造 VIMPO-style policy-implied value loss：
\begin{equation}
  \mathcal{L}_{V,\mathrm{sum}}
  = \left[
    \sum_{t\in\tau_1}\beta\Delta_t^{\mathrm{sum}}
    - \left(R(\tau_2)-b(x,\tau_1)\right)
  \right]^2.
\end{equation}

这条路线的含义是：summary-conditioned policy shift 的 cumulative score 应该对齐 suffix improvement。若 suffix 没有改善，累计 log-ratio 不应无条件变大。

\section{完整目标与推荐分阶段实现}
总目标可以写成：
\begin{equation}
  \mathcal{L}
  = \mathcal{L}_{\mathrm{suffix}}
  + \lambda_{\mathrm{prefix}}\mathcal{L}_{\mathrm{prefix}}
  + \lambda_{\mathrm{value}}\mathcal{L}_{V}
  + \lambda_{\mathrm{format}}\mathcal{L}_{\mathrm{format}}
  + \lambda_{\mathrm{kl}}\mathrm{KL}(\pi_\theta\Vert\pi_{\mathrm{anchor}}).
\end{equation}

不建议第一版同时打开所有项。推荐阶段如下：
\begin{enumerate}
  \item \textbf{SSCA-Suffix}：只优化 summary tokens 和 \path{traj2} actions。
  \item \textbf{SSCA-PrefixWeight}：加入 $\Delta_t^{\mathrm{sum}}$ 作为 \path{traj1} 的 stop-grad advantage weight。
  \item \textbf{SSCA-Value}：加入 $V_\phi(h_t,\sigma)$ 或 $\mathcal{L}_{V,\mathrm{sum}}$。
  \item \textbf{SSCA-Retrieve}：把有效 summary 写入历史经验库，形成跨 batch memory-conditioned critic。
\end{enumerate}

\section{推荐第一版算法}
\begin{lstlisting}
Algorithm: Self-Summary Policy Optimization (SSPO / SSCA)

1. Roll out tau_1 under pi_old.
2. Obtain intermediate feedback r_1.
3. Generate sigma = Summary(tau_1, r_1) under pi_old.
4. Continue rollout tau_2 under pi_old conditioned on sigma.
5. Compute suffix reward R and prefix baseline b(x, tau_1).
6. Update summary + tau_2 with A_suffix = R - b.
7. If A_suffix > 0:
     compute Delta_t^sum =
       log pi_old(a_t | h_t, sigma) - log pi_old(a_t | h_t)
     update tau_1 with stopgrad(normalized Delta_t^sum).
8. Optionally train V_phi(h_t, sigma) to predict R or A_suffix.
\end{lstlisting}

这个版本的优点是清楚、可诊断、不会过早引入复杂 critic。它把 summary 当成可训练 action，同时只在 summary 被 suffix reward 验证后才让它参与 prefix supervision。

\section{和现有工作的关系}
\begin{center}
\begin{tabularx}{\linewidth}{lXX}
\toprule
方向 & 代表工作 & 与 SSCA 的关系 \\
\midrule
Reflection / retry & Reflexion, Self-Refine, Reflect Retry Reward & 证明 reflection 可以改善后续尝试；SSCA 进一步把有效 summary 反过来监督 prefix。 \\
Policy-implied credit & VIMPO, SDPO, SC-GRPO & 提供 log-ratio / self-conditioned shift 的理论和工程先例；SSCA 将 reference 换成 summary-conditioned counterfactual policy。 \\
Agent step credit & GiGPO, GAGPO, HCAPO & 提供长程 agent step-level baseline；SSCA 不依赖 exact same-state group，而依赖 within-trajectory hindsight summary。 \\
Memory RL & MMPO, Memory-R2 & 关注 summary/memory 质量和公平 credit；SSCA 用 suffix improvement + policy shift 替代或补充 belief entropy / local rerollout。 \\
Actor-as-critic & POISE, V0.5 & 支持从 actor 内部状态或 prior 读 value；SSCA 可学习 summary-conditioned value。 \\
\bottomrule
\end{tabularx}
\end{center}

\paragraph{与 VIMPO 的关键差异。}
VIMPO 的 ratio 是 current policy 相对 frozen reference；SSCA 的 ratio 是 summary-conditioned current/old policy 相对 unconditioned current/old policy。因此 SSCA 不应被表述为“把 VIMPO 的 ref 换成 now”，更准确是：
\begin{quote}
用 summary 诱导同一 actor 的 counterfactual policy，再用 contextual log-ratio 作为 hindsight credit signal。
\end{quote}

\paragraph{与 SC-GRPO 的关键联系。}
SC-GRPO 的重要教训是：own solution 不适合直接当 distillation target，但适合当 credit filter。SSCA 应采用同样原则：
\begin{quote}
good summary is not a target; good summary is a credit filter.
\end{quote}
也就是 reward / suffix improvement 决定更新方向，summary-conditioned shift 决定哪些 prefix step/token 值得更新。

\section{实验设计}
\subsection{任务选择}
第一批任务应满足四个条件：多步 trajectory、可验证 reward、可切分 \path{traj1/traj2}、summary 确实可能改善 continuation。

\begin{center}
\begin{tabularx}{\linewidth}{lXX}
\toprule
任务 & 优点 & 风险 \\
\midrule
Countdown / math retry & 成本低，适合验证 summary RL 与 prefix score correlation & 不够 agentic \\
API / function calling & 失败原因清晰，和 Reflect Retry Reward 可比 & 环境复杂度有限 \\
WebShop & 经典 agent RL，有 GiGPO/GAGPO 对照 & state/action 模板化，summary 收益可能有限 \\
ALFWorld & 长程多步，和现有 agent credit baseline 可比 & exact state overlap 可能已经很强 \\
Code repair / unit test & \path{traj1} 失败、summary、\path{traj2} 修复很自然 & token/action 粒度对齐较难 \\
\bottomrule
\end{tabularx}
\end{center}

建议路线：
\begin{lstlisting}
cheap retry task -> validate summary RL and prefix correlation
WebShop / ALFWorld -> validate long-horizon agent benefit
\end{lstlisting}

\subsection{Ablation}
\begin{center}
\begin{tabularx}{\linewidth}{lX}
\toprule
方法 & 说明 \\
\midrule
GRPO / DAPO & trajectory-level reward baseline \\
Reflect-Retry-Reward style & 只奖励 summary/reflection tokens，不监督 prefix \\
GiGPO / GAGPO & step-level agent credit baseline \\
SSCA-Suffix & 训练 summary + \path{traj2} \\
SSCA-PrefixWeight & 加 summary-conditioned prefix advantage \\
SSCA-Value & 加 summary-conditioned value predictor \\
SSCA-NoRewardGiven & summary 不看 $r_1$，验证反馈是否必要 \\
SSCA-RandomSummary & 控制“只是增加上下文长度”的可能性 \\
SSCA-BadSummaryGateOff & suffix 未改善时关闭 prefix supervision \\
\bottomrule
\end{tabularx}
\end{center}

\subsection{诊断指标}
训练前必须先做 offline diagnostics，否则容易把错误 summary 当成监督源。

\paragraph{Summary usefulness.}
\begin{equation}
  \Delta R_{\mathrm{summary}}
  = \mathbb{E}[R(\tau_2)\mid \mathrm{with\ summary}]
  - \mathbb{E}[R(\tau_2)\mid \mathrm{generic/no\ summary}].
\end{equation}
如果该值不显著，说明 summary prompt 或任务切分本身有问题。

\paragraph{Prefix score correlation.}
\begin{equation}
  S_{\mathrm{prefix}}=\sum_t\Delta_t^{\mathrm{sum}},\qquad
  \Delta R=R(\tau_2)-b(x,\tau_1).
\end{equation}
检查 $\mathrm{corr}(S_{\mathrm{prefix}},\Delta R)$。如果相关性弱，$\Delta_t^{\mathrm{sum}}$ 不应进入训练，只能作为分析信号。

\paragraph{Key-step localization.}
在有人工或规则标注的样本上，检查高 $\Delta_t^{\mathrm{sum}}$ 是否集中在关键动作：工具选择、query reformulation、wrong branch、search/filter/compare 决策、code patch 关键行。

\paragraph{Summary faithfulness.}
检查 summary 是否引用了 \path{traj1} 中不存在的信息。可加一个轻量 verifier：
\begin{lstlisting}
summary claim must be supported by trajectory evidence
\end{lstlisting}

\subsection{建议记录的 metrics}
\begin{lstlisting}
ssca/summary_token_len_mean
ssca/summary_useful_gate_rate
ssca/suffix_adv_mean
ssca/suffix_adv_std
ssca/prefix_delta_sum_mean
ssca/prefix_delta_sum_std
ssca/prefix_delta_reward_corr
ssca/prefix_weight_entropy
ssca/value_before_mean
ssca/value_after_mean
ssca/value_delta_reward_corr
ssca/summary_faithfulness_fail_rate
ssca/random_summary_gap
\end{lstlisting}

\section{主要风险与约束}
\paragraph{Summary reward hacking.}
如果只奖励 \path{traj2}，summary 可能学会模板化文本，例如“我应该更仔细”。需要 schema、长度约束、factuality verifier 和 RandomSummary ablation。

\paragraph{Moving-reference instability.}
VIMPO 的 reference 是 frozen，SSCA 的 $\pi_{\mathrm{now}}$ 和 $\pi_{\mathrm{sum}}$ 都可能随训练漂移。第一版应使用 \path{pi_old} 计算 ratio，并对 $\Delta_t^{\mathrm{sum}}$ stop-gradient。

\paragraph{Action-dependent baseline bias.}
如果 $V_\phi(x,\tau_1,\sigma)$ 直接看当前 action 并作为 policy-gradient baseline，可能引入 bias。value predictor 第一版应作为 auxiliary loss；若用于 baseline，应采用 cross-rollout / leave-one-out 设计。

\paragraph{错误 summary 污染 prefix.}
如果 \path{traj2} 没改善，summary 不应 supervise prefix。prefix loss 必须被 suffix improvement gate 控制。

\paragraph{与 StepPPO-v3 的边界。}
当前 EXPS 主线中 StepPPO-v3 是 value-head-only 实验，不默认改变 GiGPO policy advantage、policy loss 或 rollout recipe。SSCA 应作为新的实验方向，不应写成 v3 的默认修改。

\section{预期结果}
合理预期不是第一次实验就超过所有 agent RL baseline，而是先证明三件事：
\begin{enumerate}
  \item SSCA-Suffix 比普通 GRPO 更快学会有效 reflection / summary；
  \item $\sum_t\Delta_t^{\mathrm{sum}}$ 与 suffix improvement 有正相关；
  \item 在长轨迹、routine step 多的任务上，SSCA-PrefixWeight 能提升 sample efficiency 或降低 advantage 方差。
\end{enumerate}

若上述诊断成立，再推进 SSCA-Value 和历史 summary retrieval；若诊断不成立，应优先改 summary prompt、切分点、feedback 设计，而不是加更复杂的 critic。

\section{参考入口}
\begin{itemize}
  \item VIMPO: \url{https://arxiv.org/abs/2606.20008}
  \item SC-GRPO / Learning from Own Solutions: \url{https://arxiv.org/abs/2606.18810}
  \item Self-Distillation Policy Optimization: \url{https://arxiv.org/abs/2601.20802}
  \item Reflect, Retry, Reward: \url{https://arxiv.org/abs/2505.24726}
  \item Reflexion: \url{https://arxiv.org/abs/2303.11366}
  \item Self-Refine: \url{https://arxiv.org/abs/2303.17651}
  \item STaR: \url{https://arxiv.org/abs/2203.14465}
  \item GAGPO: \url{https://arxiv.org/abs/2605.13217}
  \item Memory-R2: \url{https://arxiv.org/abs/2605.21768}
  \item MMPO: \url{https://arxiv.org/abs/2605.30159}
  \item HCAPO: \url{https://arxiv.org/abs/2603.08754}
  \item POISE: \url{https://arxiv.org/abs/2605.07579}
\end{itemize}

\end{document}
